\documentclass[justified, nobib]{tufte-handout}

\title{Notes: Euler equations \& linear hyperbolic systems}

\author[Wojciech Sadowski]{Wojciech Sadowski}

%\date{28 March 2010} % without \date command, current date is supplied

% \geometry{showframe} % display margins for debugging page layout

\usepackage{fontspec}
% \setmainfont[Renderer=Basic, Numbers=OldStyle, Scale = 1.0]{TeX Gyre Pagella}
% \setsansfont[Renderer=Basic, Scale=0.90]{TeX Gyre Heros}
% \setmonofont[Renderer=Basic]{TeX Gyre Cursor}

% Set up the spacing using fontspec features
\renewcommand\allcapsspacing[1]{{\addfontfeature{LetterSpace=15}#1}}
\renewcommand\smallcapsspacing[1]{{\addfontfeature{LetterSpace=10}#1}}

% setting ptions for included pictures
\usepackage{graphicx} % allow embedded images
\graphicspath{{graphics/}} % set of paths to search for images
\setkeys{Gin}{
  width=\linewidth,
  totalheight=\textheight,
  keepaspectratio
}

\hypersetup{colorlinks}

\usepackage[compatibility=false, font=footnotesize]{caption}
% * options are necessary to fix the bug that autoref was causing,
%   figure was constantly changed to section
\DeclareCaptionFont{blue}{\color{blue}}
\captionsetup{labelfont={blue}}

% color and color palettes
\usepackage{xcolor}

% drawings collor palette
\definecolor{p11}{HTML}{33cfff}
\definecolor{p12}{HTML}{ffa733}
\definecolor{p13}{HTML}{ff4133}


% drawing inside tex
\usepackage{tikz} 
\usetikzlibrary{decorations.markings}
\usetikzlibrary{calc}
\usetikzlibrary{shapes}

\usepackage{pgfplots}
\pgfplotsset{compat=newest}
\usepgfplotslibrary{fillbetween}
\usepackage{xcolor}

% math things
\usepackage{
  amsmath,  % extended mathematics
  amsthm,
  siunitx,  % non-stacked fractions and better unit spacing
  bm,       % bold math symbols
  physics2,  % a set of usefull symbols for differentiation 
  fixdif,
  derivative
}

\usephysicsmodule{ab, nabla.legacy}

% table things
\usepackage{
  booktabs, % book-quality tables
  multicol  % multiple column layout facilities
} 


% writing references
\usepackage{cleveref}  % must be loded after amsmath

% quoting text and friends
\usepackage{fancyvrb}         % extended verbatim environments
\fvset{fontsize=\normalsize}  % default font size for fancy-verbatim environments
\usepackage{
  dirtytalk, % provides \say command for easy quiting 
  nth       % provides \nth command for 1-st, 2-nd, etc
}        


% citing sources
\usepackage{natbib}
\setcitestyle{authoryear}
\bibliographystyle{plainnat}

% misc 
\usepackage{lipsum}   % filler text

\tikzset{
  set arrow inside/.code={\pgfqkeys{/tikz/arrow inside}{#1}},
  set arrow inside={end/.initial=>, opt/.initial=},
  /pgf/decoration/Mark/.style={
    mark/.expanded=at position #1 with
    {
      \noexpand\arrow[\pgfkeysvalueof{/tikz/arrow inside/opt}]{\pgfkeysvalueof{/tikz/arrow inside/end}}
    }
  },
  arrow inside/.style 2 args={
    set arrow inside={#1},
    postaction={
      decorate,decoration={
        markings,Mark/.list={#2}
      }
    }
  },
}

\input{../src/commands.tex}
\newtheorem{theorem}{Theorem}
\newtheorem{lemma}{Lemma}

\newcolumntype{L}[2]{>{\raggedleft#2}p{#1\textwidth}}

\input{../src/symbols.tex}



\begin{document}

\maketitle% this prints the handout title, author, and date

\tableofcontents

\section{Euler equations}
The Euler equations in conservative differential form 
(strong form):
\begin{equation}
  \begin{split}
    &\pdv{\rho}{t} + \pdv{\rho u_i}{x_i} = 0,\\
    &\pdv{\rho u_i}{t} + \pdv{\rho u_i u_j}{x_j} 
      = -\pdv{p}{x_i} + f_i,\\
    &\pdv{\rho E}{t} + \pdv{\rho E u_j}{x_j} 
      =  -\pdv{u_i p}{x_i} + \rho f_i u_i + Q
  \end{split}
  \label{eq:euler}
\end{equation}

In the integral (weak) form (the integrals taken in some
volume \(V\) with boundary \(S\)):
\begin{equation}
  \begin{split}
    &\pdv{}{t}\int\limits_V \rho \d V 
      + \oint\limits_S\rho u_i n_i \d S = 0,\\
    &\pdv{}{t}\int\limits_V \rho u_i \d V 
      + \oint\limits_S \ab( \rho u_i u_j n_j + p n_i) \d S
      = \int\limits_V f_i \d V,\\
    &\pdv{}{t}\int\limits_V \rho E \d V 
      + \oint\limits_S (\rho E + p)  u_j n_j \d S
      = \int\limits_V \ab( Q + \rho f_i u_i) \d V
  \end{split}
  \label{eq:euler-integral}
\end{equation}
We have four unknowns \(\rho, \bm{u}, E, p\) (we assume that 
\(f = Q = 0\) or are known) but we have only three equations. For an ideal
gas, we can write:
\[
  p = \rho e \ab(\gamma - 1) 
  \rightarrow 
  e = \frac{p}{\rho \ab(\gamma - 1)},
\]
where \(\gamma = c_p / c_v = 1.4\). Total energy \(E\) is equal to 
\(E = e + \bm{u}\cdot\bm{u}/2\).

Introducing the vector of conservative variables 
\(\mathcal{U}\)
\[
  \mathcal{U} = \left[
    \begin{array}{c}
      \rho\\
      \rho \bm{u}\\
      \rho E
    \end{array} \right],
\]
the advective flux \(\mathcal{F}_\mathrm{adv}\)
\[
  \mathcal{F}_\mathrm{adv} = \left[
    \begin{array}{c}
      \rho \bm{u}\\
      \rho \bm{u}\otimes\bm{u} + p \bm{I}\\
      \ab(\rho E  + p) \bm{u}
    \end{array} \right],
\]
and sources \(\mathcal{Q}\)
\[
  \mathcal{Q} = \left[
    \begin{array}{c}
      0\\
      \bm{f}\\
      Q + \rho \bm{u}\cdot\bm{f}
    \end{array} \right],
\]
we can write both in symbolic form
\[
  \pdv{\mathcal{U}}{t} + \div\mathcal{F}_\mathrm{adv} = \mathcal{Q}, 
\]
or
\[
  \pdv{}{t} \int\limits_V \mathcal{U} \d V 
  + \oint\limits_S \mathcal{F}_\mathrm{adv}\cdot \bm{n} \d S
  = \int\limits_V \mathcal{Q} \d V.
\]

\section{Euler equations in 1D}
\begin{equation}
  \begin{split}
    &\pdv{\rho}{t} + \pdv{\rho u}{x} = 0,\\
    &\pdv{\rho u}{t} + \pdv{\rho u^2}{x} 
      = -\pdv{p}{x} ,\\
    &\pdv{\rho E}{t} + \pdv{\rho E u}{x} 
      =  -\pdv{up}{x} 
  \end{split}
  \label{eq:euler1d}
\end{equation}

Introducing the vector of conservative variables 
\(\mathcal{U}\)
\[
  \mathcal{U} = \left[
    \begin{array}{c}
      \rho\\
      \rho u\\
      \rho E
    \end{array} \right],
\]
the advective flux \(\mathcal{F}_\mathrm{adv}\)
\[
  \mathcal{F}_\mathrm{adv}(\mathcal{U}) = \left[
    \begin{array}{c}
      \rho u\\
      \rho u^2 + p \\
      \ab(\rho E  + p) u
    \end{array} \right],
\]
we can write both in symbolic form
\[
  \pdv{\mathcal{U}}{t} + \pdv{\mathcal{F}_\mathrm{adv}(\mathcal{U})}{x} 
  = 0. 
\]

\section{Quasi-linear form}
\def\F{\mathcal{F}_\mathrm{adv}}
Using the chain rule, we can linearize each of the equations. The mass conservation
\[
  \begin{split}
    \pdv{\rho}{t} + \pdv{\F^\rho(\mathcal{U})}{x} 
    &= \pdv{\rho}{t} 
    + \pdv{\F^\rho(\mathcal{U})}{\mathcal{U}}\pdv{\mathcal{U}}{x}\\
    &= \pdv{\rho}{t} + \pdv{\F^\rho}{\rho}\pdv{\rho}{x} 
    + \pdv{\F^\rho}{\rho u}\pdv{\rho u}{x} 
    + \pdv{\F^\rho}{\rho E}\pdv{\rho E}{x} \\
    &= \pdv{\rho}{t} + (1)\pdv{\rho u}{x} = 0
  \end{split}
\]

\def\U{\mathcal{U}}
For momentum 
\[
  \begin{split}
    \F^{\rho u} 
    &= \rho u^2 + p = \rho u^2 + (\gamma-1)\rho e 
    =  \rho u^2 + (\gamma-1)\rho (E - u^2/2) \\
    &= \frac{\U_2^2}{\U_1} + (\gamma - 1)\ab(\U_3 - \frac{\U_2^2}{2\U_1} )
    = (3-\gamma)\frac{\U_2^2}{2\U_1}  +(\gamma - 1)\U_3\\
    &= (3-\gamma)\frac{{(\rho u)}^2}{\rho} + (\gamma - 1)\rho E   
  \end{split}
\]
so
\[
  \begin{split}
    \pdv{\rho u }{t} &+ \pdv{\F^{\rho u}(\U)}{x} 
    = \pdv{\rho u}{t} + \pdv{\F^{\rho u}(\U)}{\U}\pdv{\U}{x} \\
    &=  \pdv{\rho u}{t} 
    + \pdv{\F^{\rho u}}{\rho }\pdv{\rho }{x} 
    + \pdv{\F^{\rho u}}{\rho u}\pdv{\rho u}{x} 
    + \pdv{\F^{\rho u}}{\rho E}\pdv{\rho E}{x} \\
    &=  \pdv{\rho u}{t} 
    + \frac{ \gamma - 3 }{2} \frac{{(\rho u)}^2}{\rho^2} \pdv{\rho }{x} 
    + \ab[\frac{2\rho u }{\rho} - (\gamma -1)\frac{\rho u}{\rho} ]\pdv{\rho u}{x} + (\gamma - 1)\pdv{\rho E}{x}\\
    & = \pdv{\rho u}{t} + \frac{ \gamma - 3 }{2}u^2\pdv{\rho }{x} 
    + (3 - \gamma)u\pdv{\rho u}{x}
    + (\gamma - 1 )\pdv{\rho E}{x}
    = 0
  \end{split}
\]
For energy
\[
  \begin{split}
    \F^{\rho E} 
    &= \rho u E + p u = \rho u E + (\gamma-1)\rho e u
    = \rho u E + (\gamma-1)\rho(E - u^2/2)u \\
    &= \rho u E + (\gamma - 1 )\rho u E -(\gamma -1)\rho u^3/2
    = \gamma\rho u E -\frac{\gamma-1}{2} \rho u^3\\
    &= \gamma\frac{(\rho u)(\rho E)}{\rho} -\frac{\gamma-1}{2}\frac{{(\rho u)}^3}{\rho^2}
  \end{split}
\]
so
\[
  \begin{split}
    \pdv{\rho E }{t} &+ \pdv{\F^{\rho E}(\U)}{x} 
    = \pdv{\rho E}{t} + \pdv{\F^{\rho E}(\U)}{\U}\pdv{\U}{x} \\
    &=  \pdv{\rho E}{t} 
    + \pdv{\F^{\rho E}}{\rho }\pdv{\rho }{x} 
    + \pdv{\F^{\rho E}}{\rho u}\pdv{\rho u}{x} 
    + \pdv{\F^{\rho E}}{\rho E}\pdv{\rho E}{x} \\
    &=  \pdv{\rho E}{t} 
    + \ab[-\gamma\frac{(\rho u)(\rho E)}{\rho^2} + \ab(\gamma-1)\frac{{(\rho u)}^3}{\rho^3} ]\pdv{\rho}{x} \\
    &+ \ab[\gamma\frac{\rho E}{\rho} - \frac{3}{2}\ab(\gamma-1)\frac{{(\rho u)}^2}{\rho^2}]\pdv{\rho u}{x} 
    + \gamma\frac{\rho u}{\rho} \pdv{\rho E}{\rho} \\
    &=  \pdv{\rho E}{t} + \ab[(\gamma -1)u^3-\gamma u E ]\pdv{\rho}{x}
    + \ab[\gamma E - \frac{3}{2}(\gamma -1)u^2]\pdv{\rho u}{x}
    + \gamma u\pdv{\rho E}{\rho}
  \end{split}
\]
\def\A{\mathcal{A}}
All of this hard work, let us express our system of equations as 
\[
  \pdv{\U}{t} + \A(\U)\pdv{\U}{x} = \bm{0}
\]
or
\[
  \pdv{}{t}\left[
    \begin{array}{c}
      \rho\\
      \rho \bm{u}\\
      \rho E
    \end{array} \right] +
    \ab[\begin{array}{ccc}
      \pdv{\F^\rho}/{\rho}
      & \pdv{\F^\rho}/{\rho u}
      & \pdv{\F^\rho}/{\rho E} \\
      \pdv{\F^{\rho u}}/{\rho }
      & \pdv{\F^{\rho u}}/{\rho u}
      & \pdv{\F^{\rho u}}/{\rho E} \\
      \pdv{\F^{\rho E}}/{\rho } 
      & \pdv{\F^{\rho E}}/{\rho u}
      & \pdv{\F^{\rho E}}/{\rho E}
    \end{array}]
    \left[
   \begin{array}{c}
      \pdv{\rho}/{x}\\
      \pdv{\rho u}/{x}\\
      \pdv{\rho E}/{x}
    \end{array} \right]
  = \bm{0}.
\]
\def\A{\mathcal{A}}

The matrix \(\A\)
\[
  \A(\U) =  \left[
    \begin{array}{ccc}
      0 & 1 & 0 \\
      \dfrac{\gamma -3}{2}u^2 & (3 - \gamma)u & \gamma - 1\\
      -\gamma u E + (\gamma -1)u^3 & \gamma E - \dfrac{3(\gamma -1)}{2}u^2 & \gamma u
    \end{array} \right]
\]
is the Jacobian of the system of equations.
This is the non-conservative form of the Euler equations. 
Based on \(\A\) it can be proven that Euler equations are in fact
hyperbolic (\(\A\) has real eigenvalues and independent eigenvectors). 


\section{Primitive variables}
\def\P{\mathcal{P}}
In the similar manner we can get the primitive form, i.e., expressed 
using the primitive variables \(\P\):
\[
  \P = \left[
    \begin{array}{c}
      \rho\\
      u \\
      p
    \end{array}
  \right].
\]
\def\K{\mathcal{K}}
We can simply make a change of variables using the chain rule, 
\[
  \pdv{\U}{t} = \pdv{\U}{\P}\pdv{\P}{t} = 
  \ab[
    \begin{array}{ccc}
      \pdv{\rho  }/{\rho} & \pdv{\rho  }/{u} & \pdv{\rho  }/{p}\\
      \pdv{\rho u}/{\rho} & \pdv{\rho u}/{u} & \pdv{\rho u}/{p}\\
      \pdv{\rho E}/{\rho} & \pdv{\rho E}/{u} & \pdv{\rho E}/{p}
    \end{array}
  ] \pdv{\P}{t} = \K\pdv{\P}{t}
\] 
We only need to clarify the transformation of
energy into pressure
\[
  \rho E = \rho e + \rho u^2/2 = \frac{p}{\gamma -1 } + \rho u^2/2,
\]
From that we have
\[
  \pdv{\rho E}{\rho} = \pdv{}{\rho}\ab(\frac{p}{\gamma -1 } + \rho u^2/2)
  = \frac{u^2}{2},
\]
\[
  \pdv{\rho E}{u} = \pdv{}{u}\ab(\frac{p}{\gamma -1 } + \rho u^2/2)
  = \rho u,
\]
and finally
\[
  \pdv{\rho E}{p} = \pdv{}{p}\ab(\frac{p}{\gamma -1 } + \rho u^2/2)
  = \frac{1}{\gamma -1}.
\]

The Jacobian of variable change \(\pdv{\U}/{\P}\) can be easily computed
\[
  \pdv{\U}{\P} = 
  \ab[
    \begin{array}{ccc}
      1 & 0 & 0\\
      u & \rho & 0\\
      u^2/2 & \rho u & 1/(\gamma -1)
    \end{array}
  ]
\]
and inserted into our original equation:
\[
  \K\pdv{\P}{t} + \A(\U)\K\pdv{\P}{x} = \bm{0}.
\]
Is \(\K\) invertible? It is a lower triangular matrix so
\[
  \det \K = \frac{\rho}{\gamma - 1} > 0
\]
and 
\[
  \K^{-1} = \ab[
    \begin{array}{ccc}
      1 & 0 & 0\\
      -u/\rho & 1/\rho & 0\\
      (\gamma -1)u^2/2 & u(1 - \gamma)& (\gamma -1)
    \end{array}
  ] 
\]
We know that \(\K^{-1}\K = \bm{I}\), so we can multiply the equations 
from left side by \(\K^{-1}\) and get the following form:
\[
  \pdv{\P}{t} + \K^{-1}\A(\U)\K\pdv{\P}{x} = \bm{0}.
\]
The matrix \(\A(\P) = \K^{-1}\A(\U)\K\) is given as\footnote{Check that and send it to me!}
\[
  \A(\P) = \ab[
    \begin{array}{ccc}
      u & \rho & 0\\
      0 & u & 1/\rho\\
      0 & \rho a^2 & u
    \end{array}
  ],
\]
where \(a = \sqrt{\gamma p /\rho}\) is the speed of sound. The 
matrix form can be expanded to 
\[
  \begin{split} 
    \pdv{\rho}{t} &+ u\pdv{\rho}{x} + \rho\pdv{u}{x} = 0\\
    \pdv{u}{t}    &+ u\pdv{u}{x} + \frac{1}{\rho}\pdv{p}{x} = 0 \\
    \pdv{p}{t}    &+ \rho a^2 \pdv{u}{x} + u\pdv{p}{x} = 0.
  \end{split}
\]
\subsection{Entropy equation}
If we consider entropy \(S = p/\rho^\gamma\) we can compute 
\[
  \pdv{S}{t} = \frac{1}{\rho^\gamma}\pdv{p}{t} + \pdv{\rho^{-\gamma}}{t}p 
  = \frac{1}{\rho^\gamma}\pdv{p}{t} - \gamma p \rho^{-\gamma - 1}\pdv{\rho}{t}
  = \frac{1}{\rho^\gamma}\ab(\pdv{p}{t} - \frac{\gamma p}{\rho}\pdv{\rho}{t})
\]
which can be further simplified using the speed of sound and \nth{1} and
\nth{3} equation in the primitive form and substitute the substantial
derivatives:
\[
  \begin{split}
    \frac{1}{\rho^\gamma}\ab(\pdv{p}{t} - a^2\pdv{\rho}{t})
    & = \frac{1}{\rho^\gamma}\ab[-\rho a^2 \pdv{u}{x} - u\pdv{p}{x} + a^2\ab(u\pdv{\rho}{x} + \rho\pdv{u}{x}) ]\\
    & = \frac{-u}{\rho^\gamma}\ab( \pdv{p}{x} - a^2\pdv{\rho}{x} ) = -u\pdv{S}{x}
  \end{split}
\]
Equating one to the other gives us the equation for entropy that is valid in 
smooth regions of the flow 
\[
  \pdv{S}{t} + u\pdv{S}{x} = \mdv{S}{t}=0,
\]
which implies that entropy remains constant along the flow paths.

\section{General linear system}
A general linear system is given as 
\[
  \pdv{\bm{u}}{t} + \bm{A}\pdv{\bm{u}}{x} = \bm{0}, \quad \bm{u} = {\ab[u_1 \dots u_n]}^T, 
\]
where the flux Jacobian \(\bm{A}\) is constant, has \(n\) distinct real
eigenvalues \(\lambda_i\) and \(n\) linearly independent eigenvectors
\(\bm{k}_i\).

\subsection{Diagonalization}
\newcommand{\vb}[1]{\bm{#1}} 
If we solve the eigenvalue problem \(\bm{A}\bm{k}
-\lambda\bm{k} =\bm{0}\), we can create a matrix \(\vb{K}=\ab[\vb{k}_1 \dots
\vb{k}_n]\). Then the coefficient matrix \(\vb{A}\) can be expressed as
\[
  \bm{A} = \bm{K}\bm{\Lambda}\bm{K}^{-1},\quad 
  \bm{\Lambda} = \ab[
    \begin{array}{ccc} 
      \lambda_1 &        & \\
               & \ddots & \\
               &        &  \lambda_n
    \end{array}].
\]
\subsection{Characteristic variables}
We can also introduce new set of variables \(\vb{w} = \bm{K}\vb{u}\), and 
substitute it into our equation leading to
\[
  \pdv{\bm{w}}{t} + \bm{\Lambda}\pdv{\bm{w}}{x} = \bm{0}.
\]
This is a decoupled system!\footnote{\(\bm\Lambda\) is diagonal.}
We also see, that \(\bm{w}_i\) is the coefficient of \(i\)-th eigenvector 
when \(\bm{u}\) is represented in the basis of eigenvectors
\[
  \bm{u}(x,t) = \bm{K}\bm{w} = \ab[\vb{k}_1 \dots \vb{k}_n]{\ab[w_1 \dots w_n]}^T=
  \sum\limits_i^n w_i(x,t)\vb{k}_i
\]
If we know the initial condition to \(\vb{u} = \vb{u}^{(0)}\), we can transform it
as well, and then solve each of the equations separately using the 
method of characteristics:
\[
  \bm{u}(x,t) = \sum\limits_i^n w^{(0)}_i(x-\lambda_i t)\vb{k}_i.
\]


\end{document}
